\section{Three-dimensional Weyl and Dirac}
\label{p:weyl-dirac-3d}

The transition to three dimensions adds two essential features absent in
$2{+}1$: a real chirality assignment to each band-crossing node, and a
controlled \emph{splitting} between two opposite-chirality Weyl nodes
that together form a Dirac fermion.

\subsection{Two-band Weyl and four-band Dirac}
The two-band massless Weyl Hamiltonian is $\HW = v\,\vsigma\cdot\vp$. The
Pauli algebra gives $(\vsigma\cdot\vp)^{2} = \lvert\vp\rvert^{2}\,\bbI$,
hence a linear spectrum $E_{\pm}=\pm v\lvert\vp\rvert$ and no local mass
term: no nonzero element of $\{\bbI,\sigma_{i}\}$ anticommutes with all
three $\sigma_{i}$ simultaneously, so a Weyl node cannot be gapped by an
on-site perturbation (in the sense relevant here)~\cite{Burkov2018}. The
chirality of the node is read directly from the Jacobian
$J_{ab} = \partial h_{a}/\partial k_{b}$ evaluated at the node, namely
$\chi = \operatorname{sgn}\det J$. For $\HW = v\,\vsigma\cdot\vp$,
$J=v\,\bbI_{3}$, so $\chi = \operatorname{sgn} v$; the parity-related node
carries opposite chirality
(\texttt{validators/chirality.py}).

A four-component Dirac fermion is then built by combining two
opposite-chirality Weyl blocks through a mass term that anticommutes
with the kinetic part:
\begin{equation}
\HD = v\,\tau_{z}\otimes(\vsigma\cdot\vp) + m\,\tau_{x}\otimes\bbI ,
\qquad
\bigl\{\tau_{z}\!\otimes\!(\vsigma\!\cdot\!\vp),\,\tau_{x}\!\otimes\!\bbI\bigr\} = 0 .
\end{equation}
Squaring gives the canonical relativistic form
$\HD^{2} = (v^{2}\lvert\vp\rvert^{2}+m^{2})\,\bbI_{4}$, hence a gap $2m$
at $\vp=0$ (\texttt{validators/estabilidad\_dirac.py}). When $m\to 0$ the
four-band Dirac fermion separates into two two-band Weyl fermions of
opposite chirality, which is exactly the structure Nielsen--Ninomiya
requires of any lattice realization.

% Figure 2: Weyl splitting under axial perturbation
\begin{figure}[t]
\centering
\begin{tikzpicture}[
  every node/.style={font=\small},
  axis/.style={->, >=Latex, thick},
  cone/.style={very thick},
  band/.style={very thick, color=#1},
  helpline/.style={dashed, color=black!40},
]
% --- Left: single Dirac point at p_z = 0 (b = 0)
\begin{scope}[shift={(0,0)}]
  \draw[axis] (-2,0) -- (2,0) node[right] {\(p_z\)};
  \draw[axis] (0,-2) -- (0,2) node[above] {\(E\)};
  % cone at origin
  \draw[band=red!70!black] (-1.8,1.8) -- (1.8,-1.8);
  \draw[band=blue!70!black] (-1.8,-1.8) -- (1.8,1.8);
  \node at (0,-2.6) {\textbf{(a)} $b = 0$: single Dirac point};
\end{scope}

% --- Right: two Weyl points at p_z = +/- b/v
\begin{scope}[shift={(7,0)}]
  \draw[axis] (-3,0) -- (3,0) node[right] {\(p_z\)};
  \draw[axis] (0,-2) -- (0,2) node[above] {\(E\)};
  % left Weyl node at p_z = -b/v (call it -1.2)
  \draw[band=red!70!black] (-2.6,1.4) -- (-1.2,0) -- (0.2,1.4);
  \draw[band=blue!70!black] (-2.6,-1.4) -- (-1.2,0) -- (0.2,-1.4);
  % right Weyl node at p_z = +b/v (call it 1.2)
  \draw[band=red!70!black] (-0.2,1.4) -- (1.2,0) -- (2.6,1.4);
  \draw[band=blue!70!black] (-0.2,-1.4) -- (1.2,0) -- (2.6,-1.4);
  % markers
  \draw[helpline] (-1.2,-2) -- (-1.2,2);
  \draw[helpline] (1.2,-2) -- (1.2,2);
  \node[below] at (-1.2,-0.1) {\scriptsize \(-b/v\)};
  \node[below] at (1.2,-0.1) {\scriptsize \(+b/v\)};
  \node[above right, font=\scriptsize, red!70!black] at (1.4,0.6) {\(\chi=+1\)};
  \node[above left,  font=\scriptsize, blue!70!black] at (-1.4,0.6) {\(\chi=-1\)};
  \node at (0,-2.6) {\textbf{(b)} \(b \ne 0\): two Weyl nodes split by \(2b/v\)};
\end{scope}
\end{tikzpicture}
\caption{Controlled splitting of a Dirac point into two Weyl nodes by an
axial Zeeman-like perturbation $b\,\bbI\otimes\sigma_z$. The single Dirac
crossing at $\vec p = 0$ in (a) separates into two Weyl crossings of
opposite chirality along $p_z$ in (b), at $p_z = \pm b/v$.}
\label{p:fig:splitting}
\end{figure}


\subsection{Controlled splitting Dirac \texorpdfstring{$\to$}{->} two Weyls}
The cleanest way to obtain a separated Weyl pair from a Dirac point is
to add a Zeeman-like axial term aligned with one spatial direction:
\begin{equation}
\Hop \;=\; v\,\tau_{z}\otimes(\vsigma\cdot\vp)
       \;+\; b\,\bbI\otimes\sigma_{z} .
\end{equation}
Squaring,
\begin{equation}
\Hop^{2}\;=\;(v^{2}\lvert\vp\rvert^{2} + b^{2})\,\bbI_{4}
   \;+\; 2\,v\,b\,p_{z}\,(\tau_{z}\otimes\bbI) ,
\end{equation}
which is block-diagonal in $\tau_{z}$. The two blocks have zero modes at
$\vp=(0,0,\mp b/v)$ respectively, so the original Dirac node splits into
two Weyl nodes along $p_{z}$ at separation $2b/v$, as illustrated in
Fig.~\ref{p:fig:splitting} (\texttt{validators/desdoblamiento.py}).

\subsection{Lattice realization with Wilson term}
A concrete three-dimensional lattice realization that we use as input to
the protospace construction of Sec.~\ref{p:open-problems} is the
Wilson--Dirac Hamiltonian
\begin{equation}
\Hop(\vk) \;=\; v \sum_{i}\sin(k_{i}a)\,\alpha_{i}
       + \!\Bigl(m + r\!\sum_{i}\bigl(1-\cos(k_{i}a)\bigr)\Bigr)\,\beta ,
\label{p:eq:Wilson3D}
\end{equation}
with $\alpha_{i}=\tau_{x}\otimes\sigma_{i}$ and $\beta=\tau_{z}\otimes\bbI$.
The Wilson term~\cite{Wilson1974} lifts the seven naive doublers of the
cubic Brillouin zone, leaving a single light Dirac fermion at $\vk=0$;
the doublers acquire masses $m+2r,m+4r,m+6r$ depending on how many of
their momentum components sit at the corner $k_{i}=\pi/a$
(\texttt{validators/wilson\_subsector.py},
\texttt{validators/protoespacio\_global.py}). This is the lattice
analogue of the level hierarchy of Sec.~\ref{p:setup}: a microscopic
exact Hamiltonian at Level~I whose infrared sector at Level~II is the
canonical four-component Dirac fermion, organized at Level~III by the
Clifford algebra of Sec.~\ref{p:covariant}.
